\subsection{Partition basics}

\subsubsection{Definitions}

The following definition is built in analogy to Definition \ref{def.fps.comps}:

\begin{definition}
\label{def.pars.parts}\textbf{(a)} An \emph{(integer) partition} means a
(finite) weakly decreasing tuple of positive integers -- i.e., a finite tuple
$\left(  \lambda_{1},\lambda_{2},\ldots,\lambda_{m}\right)  $ of positive
integers such that $\lambda_{1}\geq\lambda_{2}\geq\cdots\geq\lambda_{m}$.

Thus, partitions are the same as weakly decreasing compositions. Hence, the
notions of \emph{size} and \emph{length} of a partition are automatically
defined, since we have defined them for compositions (in Definition
\ref{def.fps.comps}). \medskip

\textbf{(b)} The \emph{parts} of a partition $\left(  \lambda_{1},\lambda
_{2},\ldots,\lambda_{m}\right)  $ are simply its entries $\lambda_{1}%
,\lambda_{2},\ldots,\lambda_{m}$. \medskip

\textbf{(c)} Let $n\in\mathbb{Z}$. A \emph{partition of }$n$ means a partition
whose size is $n$. \medskip

\textbf{(d)} Let $n\in\mathbb{Z}$ and $k\in\mathbb{N}$. A \emph{partition of
}$n$\emph{ into }$k$\emph{ parts} is a partition whose size is $n$ and whose
length is $k$.
\end{definition}

\begin{example}
\label{exa.pars.pars5}The partitions of $5$ are%
\[
\left(  5\right)  ,\ \ \left(  4,1\right)  ,\ \ \left(  3,2\right)
,\ \ \left(  3,1,1\right)  ,\ \ \left(  2,2,1\right)  ,\ \ \left(
2,1,1,1\right)  ,\ \ \left(  1,1,1,1,1\right)  .
\]

\end{example}

\begin{definition}
\label{def.pars.pn-pkn}\textbf{(a)} Let $n\in\mathbb{Z}$ and $k\in\mathbb{N}$.
Then, we set%
\[
p_{k}\left(  n\right)  :=\left(  \text{\# of partitions of }n\text{ into
}k\text{ parts}\right)  .
\]


\textbf{(b)} Let $n\in\mathbb{Z}$. Then, we set%
\[
p\left(  n\right)  :=\left(  \text{\# of partitions of }n\right)  .
\]
This is called the $n$\emph{-th partition number}.
\end{definition}

\begin{example}
Our above list of partitions of $5$ reveals that%
\begin{align*}
p_{0}\left(  5\right)   &  =0;\\
p_{1}\left(  5\right)   &  =1;\\
p_{2}\left(  5\right)   &  =2;\\
p_{3}\left(  5\right)   &  =2;\\
p_{4}\left(  5\right)   &  =1;\\
p_{5}\left(  5\right)   &  =1;\\
p_{k}\left(  5\right)   &  =0\ \ \ \ \ \ \ \ \ \ \text{for any }k>5;
\end{align*}
and finally $p\left(  5\right)  =7$.
\end{example}

Here are the values of $p\left(  n\right)  $ for the first $15$ nonnegative
integers $n$:%
\[%
\begin{tabular}
[c]{|c||c|c|c|c|c|c|c|c|c|c|c|c|c|c|c|}\hline
$n$ & $0$ & $1$ & $2$ & $3$ & $4$ & $5$ & $6$ & $7$ & $8$ & $9$ & $10$ & $11$
& $12$ & $13$ & $14$\\\hline
$p\left(  n\right)  $ & $1$ & $1$ & $2$ & $3$ & $5$ & $7$ & $11$ & $15$ & $22$
& $30$ & $42$ & $56$ & $77$ & $101$ & $135$\\\hline
\end{tabular}
\ \ \ \ .
\]
The sequence $\left(  p\left(  0\right)  ,p\left(  1\right)  ,p\left(
2\right)  ,\ldots\right)  $ is remarkable for being an integer sequence that
grows faster than polynomially, but still considerably slower than
exponentially. (See (\ref{eq.pars.asymptotic-pn}) for an asymptotic
expansion.) This not-too-fast growth (for instance, $p\left(  100\right)
=190\ 569\ 292$ is far smaller than $2^{100}$) makes integer partitions rather
convenient for computer experiments.

\subsubsection{Simple properties of partition numbers}

We will next state some elementary properties of $p_{k}\left(  n\right)  $ and
$p\left(  n\right)  $, but first we introduce a few very basic notations:

\begin{definition}
\label{def.pars.iverson}We will use the \emph{Iverson bracket notation}: If
$\mathcal{A}$ is a logical statement, then $\left[  \mathcal{A}\right]  $
means the \emph{truth value} of $\mathcal{A}$; this is the integer $%
\begin{cases}
1, & \text{if }\mathcal{A}\text{ is true};\\
0, & \text{if }\mathcal{A}\text{ is false}.
\end{cases}
$
\end{definition}

For example, $\left[  2+2=4\right]  =1$ and $\left[  2+2=5\right]  =0$.

Note that the Kronecker delta notation is a particular case of the Iverson
bracket: We have
\[
\delta_{i,j}=\left[  i=j\right]  \ \ \ \ \ \ \ \ \ \ \text{for any objects
}i\text{ and }j.
\]


\begin{definition}
\label{def.pars.floor-ceil}Let $a$ be a real number.

Then, $\left\lfloor a\right\rfloor $ (called the \emph{floor} of $a$) means
the largest integer that is $\leq a$.

Likewise, $\left\lceil a\right\rceil $ (called the \emph{ceiling} of $a$)
means the smallest integer that is $\geq a$.
\end{definition}

For example, the number $\pi\approx3.14$ satisfies $\left\lfloor
\pi\right\rfloor =3$ and $\left\lceil \pi\right\rceil =4$ and $\left\lfloor
-\pi\right\rfloor =-4$ and $\left\lceil -\pi\right\rceil =-3$. For another
example, $\left\lfloor n\right\rfloor =\left\lceil n\right\rceil =n$ for each
$n\in\mathbb{Z}$.

The following proposition collects various basic properties of the numbers
introduced in Definition \ref{def.pars.pn-pkn}:

\begin{proposition}
\label{prop.pars.basics}Let $n\in\mathbb{Z}$ and $k\in\mathbb{N}$. \medskip

\textbf{(a)} We have $p_{k}\left(  n\right)  =0$ whenever $n<0$ and
$k\in\mathbb{N}$. \medskip

\textbf{(b)} We have $p_{k}\left(  n\right)  =0$ whenever $k>n$. \medskip

\textbf{(c)} We have $p_{0}\left(  n\right)  =\left[  n=0\right]  $. \medskip

\textbf{(d)} We have $p_{1}\left(  n\right)  =\left[  n>0\right]  $. \medskip

\textbf{(e)} We have $p_{k}\left(  n\right)  =p_{k}\left(  n-k\right)
+p_{k-1}\left(  n-1\right)  $ whenever $k>0$. \medskip

\textbf{(f)} We have $p_{2}\left(  n\right)  =\left\lfloor n/2\right\rfloor $
whenever $n\in\mathbb{N}$. \medskip

\textbf{(g)} We have $p\left(  n\right)  =p_{0}\left(  n\right)  +p_{1}\left(
n\right)  +\cdots+p_{n}\left(  n\right)  $ whenever $n\in\mathbb{N}$. \medskip

\textbf{(h)} We have $p\left(  n\right)  =0$ whenever $n<0$.
\end{proposition}

\begin{proof}
[Proof of Proposition \ref{prop.pars.basics} (sketched).]\textbf{(a)} The size
of a partition is always nonnegative (being a sum of positive integers). Thus,
a negative number $n$ has no partitions whatsoever. Thus, $p_{k}\left(
n\right)  =0$ whenever $n<0$ and $k\in\mathbb{N}$. \medskip

\textbf{(b)} If $\left(  \lambda_{1},\lambda_{2},\ldots,\lambda_{k}\right)  $
is a partition, then $\lambda_{i}\geq1$ for each $i\in\left\{  1,2,\ldots
,k\right\}  $ (because a partition is a tuple of positive integers, i.e., of
integers $\geq1$). Hence, if $\left(  \lambda_{1},\lambda_{2},\ldots
,\lambda_{k}\right)  $ is a partition of $n$ into $k$ parts, then%
\begin{align*}
n  &  =\lambda_{1}+\lambda_{2}+\cdots+\lambda_{k}\ \ \ \ \ \ \ \ \ \ \left(
\text{since }\left(  \lambda_{1},\lambda_{2},\ldots,\lambda_{k}\right)  \text{
is a partition of }n\right) \\
&  \geq\underbrace{1+1+\cdots+1}_{k\text{ times}}\ \ \ \ \ \ \ \ \ \ \left(
\text{since }\lambda_{i}\geq1\text{ for each }i\in\left\{  1,2,\ldots
,k\right\}  \right) \\
&  =k.
\end{align*}
Thus, a partition of $n$ into $k$ parts cannot satisfy $k>n$. Thus, no such
partitions exist if $k>n$. In other words, $p_{k}\left(  n\right)  =0$ if
$k>n$. \medskip

\textbf{(c)} The integer $0$ has a unique partition into $0$ parts, namely the
empty tuple $\left(  {}\right)  $. A nonzero integer $n$ cannot have any
partitions into $0$ parts, since the empty tuple has size $0\neq n$. Thus,
$p_{0}\left(  n\right)  $ equals $1$ for $n=0$ and equals $0$ for $n\neq0$. In
other words, $p_{0}\left(  n\right)  =\left[  n=0\right]  $. \medskip

\textbf{(d)} Any positive integer $n$ has a unique partition into $1$ part --
namely, the $1$-tuple $\left(  n\right)  $. On the other hand, if $n$ is not
positive, then this $1$-tuple is not a partition, so in this case $n$ has no
partition into $1$ part. Thus, $p_{1}\left(  n\right)  $ equals $1$ if $n$ is
positive and equals $0$ otherwise. In other words, $p_{1}\left(  n\right)
=\left[  n>0\right]  $. \medskip

\textbf{(e)} Assume that $k>0$. We must prove that $p_{k}\left(  n\right)
=p_{k}\left(  n-k\right)  +p_{k-1}\left(  n-1\right)  $.

We consider all partitions of $n$ into $k$ parts. We classify these partitions
into two types:

\begin{itemize}
\item \emph{Type 1} consists of all partitions that have $1$ as a part.

\item \emph{Type 2} consists of all partitions that don't.
\end{itemize}

For example, here are the partitions of $5$ along with their types:%
\[
\underbrace{\left(  4,1\right)  ,\ \ \left(  3,1,1\right)  ,\ \ \left(
2,2,1\right)  ,\ \ \left(  2,1,1,1\right)  ,\ \ \left(  1,1,1,1,1\right)
}_{\text{Type 1}},\ \ \underbrace{\left(  5\right)  ,\ \ \left(  3,2\right)
}_{\text{Type 2}}.
\]


Let us count the type-1 partitions and the type-2 partitions separately.

Any type-1 partition has $1$ as a part, therefore as its last part (because it
is weakly decreasing). Hence, any type-1 partition has the form $\left(
\lambda_{1},\lambda_{2},\ldots,\lambda_{k-1},1\right)  $. If $\left(
\lambda_{1},\lambda_{2},\ldots,\lambda_{k-1},1\right)  $ is a type-1 partition
(of $n$ into $k$ parts), then $\left(  \lambda_{1},\lambda_{2},\ldots
,\lambda_{k-1}\right)  $ is a partition of $n-1$ into $k-1$ parts. Thus, we
have a map%
\begin{align*}
\left\{  \text{type-1 partitions of }n\text{ into }k\text{ parts}\right\}   &
\rightarrow\left\{  \text{partitions of }n-1\text{ into }k-1\text{
parts}\right\}  ,\\
\left(  \lambda_{1},\lambda_{2},\ldots,\lambda_{k-1},1\right)   &
\mapsto\left(  \lambda_{1},\lambda_{2},\ldots,\lambda_{k-1}\right)  .
\end{align*}
This map is a bijection (since it has an inverse map, which simply inserts a
$1$ at the end of a partition). Thus, the bijection principle shows that%
\begin{align*}
&  \left(  \text{\# of type-1 partitions of }n\text{ into }k\text{
parts}\right) \\
&  =\left(  \text{\# of partitions of }n-1\text{ into }k-1\text{
parts}\right)  =p_{k-1}\left(  n-1\right)
\end{align*}
(by the definition of $p_{k-1}\left(  n-1\right)  $).

Now let us count type-2 partitions. A type-2 partition does not have $1$ as a
part; hence, all its parts are larger than $1$ (because all its parts are
positive integers), and therefore we can subtract $1$ from each part and still
have a partition in front of us. To be more specific: If $\left(  \lambda
_{1},\lambda_{2},\ldots,\lambda_{k}\right)  $ is a type-2 partition of $n$
into $k$ parts, then subtracting $1$ from each of its parts produces the
$k$-tuple $\left(  \lambda_{1}-1,\lambda_{2}-1,\ldots,\lambda_{k}-1\right)  $,
which is a partition of $n-k$ into $k$ parts. Hence, we have a map%
\begin{align*}
\left\{  \text{type-2 partitions of }n\text{ into }k\text{ parts}\right\}   &
\rightarrow\left\{  \text{partitions of }n-k\text{ into }k\text{
parts}\right\}  ,\\
\left(  \lambda_{1},\lambda_{2},\ldots,\lambda_{k}\right)   &  \mapsto\left(
\lambda_{1}-1,\lambda_{2}-1,\ldots,\lambda_{k}-1\right)  .
\end{align*}
This map is a bijection (since it has an inverse map, which simply adds $1$ to
each entry of a partition). Thus, the bijection principle shows that%
\begin{align*}
&  \left(  \text{\# of type-2 partitions of }n\text{ into }k\text{
parts}\right) \\
&  =\left(  \text{\# of partitions of }n-k\text{ into }k\text{ parts}\right)
=p_{k}\left(  n-k\right)
\end{align*}
(by the definition of $p_{k}\left(  n-k\right)  $).

Since any partition of $n$ into $k$ parts is either type-1 or type-2 (but not
both at the same time), we now have%
\begin{align*}
&  \left(  \text{\# of partitions of }n\text{ into }k\text{ parts}\right) \\
&  =\underbrace{\left(  \text{\# of type-1 partitions of }n\text{ into
}k\text{ parts}\right)  }_{=p_{k-1}\left(  n-1\right)  }+\underbrace{\left(
\text{\# of type-2 partitions of }n\text{ into }k\text{ parts}\right)
}_{=p_{k}\left(  n-k\right)  }\\
&  =p_{k-1}\left(  n-1\right)  +p_{k}\left(  n-k\right)  =p_{k}\left(
n-k\right)  +p_{k-1}\left(  n-1\right)  .
\end{align*}
Since the left hand side of this equality is $p_{k}\left(  n\right)  $, we
thus have proved that $p_{k}\left(  n\right)  =p_{k}\left(  n-k\right)
+p_{k-1}\left(  n-1\right)  $. \medskip

\textbf{(f)} Let $n\in\mathbb{N}$. The partitions of $n$ into $2$ parts are%
\[
\left(  n-1,1\right)  ,\ \ \left(  n-2,2\right)  ,\ \ \left(  n-3,3\right)
,\ \ \ldots,\ \ \left(  \underbrace{n-\left\lfloor n/2\right\rfloor
}_{=\left\lceil n/2\right\rceil },\left\lfloor n/2\right\rfloor \right)  .
\]
Thus there are $\left\lfloor n/2\right\rfloor $ of them. In other words,
$p_{2}\left(  n\right)  =\left\lfloor n/2\right\rfloor $. \medskip

\textbf{(g)} Let $n\in\mathbb{N}$. Any partition of $n$ must have $k$ parts
for some $k\in\mathbb{N}$. Thus,%
\begin{align*}
p\left(  n\right)   &  =\sum_{k\in\mathbb{N}}p_{k}\left(  n\right)
=\sum_{k=0}^{n}p_{k}\left(  n\right)  +\sum_{k=n+1}^{\infty}\underbrace{p_{k}%
\left(  n\right)  }_{\substack{=0\\\text{(by Proposition
\ref{prop.pars.basics} \textbf{(b)})}}}\\
&  =\sum_{k=0}^{n}p_{k}\left(  n\right)  =p_{0}\left(  n\right)  +p_{1}\left(
n\right)  +\cdots+p_{n}\left(  n\right)  .
\end{align*}


\textbf{(h)} Same argument as for \textbf{(a)}.
\end{proof}

\subsubsection{The generating function}

Proposition \ref{prop.pars.basics} \textbf{(e)} is a recursive formula that
makes it not too hard to compute $p_{k}\left(  n\right)  $ for reasonably
small values of $n$ and $k$. Then, using Proposition \ref{prop.pars.basics}
\textbf{(g)}, we can compute $p\left(  n\right)  $ from these $p_{k}\left(
n\right)  $'s. However, one might want a better, faster method.

To get there, let me first express the generating function of the numbers
$p\left(  n\right)  $:

\begin{theorem}
\label{thm.pars.main-gf}In the FPS ring $\mathbb{Z}\left[  \left[  x\right]
\right]  $, we have%
\[
\sum_{n\in\mathbb{N}}p\left(  n\right)  x^{n}=\prod_{k=1}^{\infty}\dfrac
{1}{1-x^{k}}.
\]


(The product on the right hand side is well-defined, since multiplying a FPS
by $\dfrac{1}{1-x^{k}}$ does not affect its first $k$ coefficients.)
\end{theorem}

\begin{example}
Let us check the above equality \textquotedblleft up to $x^{5}$%
\textquotedblright, i.e., let us compare the coefficients of $x^{i}$ for
$i<5$. (In doing so, we can ignore all powers of $x$ higher than $x^{4}$.) We
have%
\begin{align*}
\prod_{k=1}^{\infty}\dfrac{1}{1-x^{k}}  &  =\dfrac{1}{1-x^{1}}\cdot\dfrac
{1}{1-x^{2}}\cdot\dfrac{1}{1-x^{3}}\cdot\dfrac{1}{1-x^{4}}\cdot\cdots\\
&  =\ \ \ \left(  1+x+x^{2}+x^{3}+x^{4}+\cdots\right) \\
&  \ \ \ \ \ \cdot\left(  1+x^{2}+x^{4}+\cdots\right) \\
&  \ \ \ \ \ \cdot\left(  1+x^{3}+\cdots\right) \\
&  \ \ \ \ \ \cdot\left(  1+x^{4}+\cdots\right) \\
&  \ \ \ \ \ \cdot\left(  1+\cdots\right) \\
&  \ \ \ \ \ \cdot\left(  1+\cdots\right) \\
&  \ \ \ \ \ \cdot\cdots\\
&  =1+x+2x^{2}+3x^{3}+5x^{4}+\cdots\\
&  =p\left(  0\right)  +p\left(  1\right)  x+p\left(  2\right)  x^{2}+p\left(
3\right)  x^{3}+p\left(  4\right)  x^{4}+\cdots.
\end{align*}

\end{example}

\begin{proof}
[Proof of Theorem \ref{thm.pars.main-gf}.]We have%
\begin{align*}
&  \prod_{k=1}^{\infty}\underbrace{\dfrac{1}{1-x^{k}}}_{=1+x^{k}+x^{2k}%
+x^{3k}+\cdots}\\
&  =\prod_{k=1}^{\infty}\underbrace{\left(  1+x^{k}+x^{2k}+x^{3k}%
+\cdots\right)  }_{=\sum\limits_{u\in\mathbb{N}}x^{ku}}=\prod_{k=1}^{\infty
}\ \ \sum_{u\in\mathbb{N}}x^{ku}\\
&  =\sum_{\substack{\left(  u_{1},u_{2},u_{3},\ldots\right)  \in
\mathbb{N}^{\infty}\text{ is}\\\text{essentially finite}}}x^{1u_{1}}x^{2u_{2}%
}x^{3u_{3}}\cdots\ \ \ \ \ \ \ \ \ \ \left(
\begin{array}
[c]{c}%
\text{here, we expanded the product}\\
\text{using Proposition \ref{prop.fps.prodrule-inf-infN}}%
\end{array}
\right) \\
&  =\sum_{\substack{\left(  u_{1},u_{2},u_{3},\ldots\right)  \in
\mathbb{N}^{\infty}\text{ is}\\\text{essentially finite}}}x^{1u_{1}%
+2u_{2}+3u_{3}+\cdots}\\
&  =\sum_{n\in\mathbb{N}}\left\vert Q_{n}\right\vert x^{n},
\end{align*}
where
\[
Q_{n}=\left\{  \left(  u_{1},u_{2},u_{3},\ldots\right)  \in\mathbb{N}^{\infty
}\text{ essentially finite }\mid\ 1u_{1}+2u_{2}+3u_{3}+\cdots=n\right\}  .
\]


Thus, it will suffice to show that%
\[
\left\vert Q_{n}\right\vert =p\left(  n\right)  \ \ \ \ \ \ \ \ \ \ \text{for
each }n\in\mathbb{N}.
\]


Let us fix $n\in\mathbb{N}$. We want to construct a bijection from $Q_{n}$ to
$\left\{  \text{partitions of }n\right\}  $.

Here is how to do this: For any $\left(  u_{1},u_{2},u_{3},\ldots\right)  \in
Q_{n}$, define a partition%
\begin{align*}
\pi\left(  u_{1},u_{2},u_{3},\ldots\right)  :=  &  \left(  \text{the partition
that contains each }i\text{ exactly }u_{i}\text{ times}\right) \\
=  &  \left(  \ldots,\underbrace{3,3,\ldots,3}_{u_{3}\text{ times}%
},\underbrace{2,2,\ldots,2}_{u_{2}\text{ times}},\underbrace{1,1,\ldots
,1}_{u_{1}\text{ times}}\right)  .
\end{align*}
This is a partition of $n$, since its size is $1u_{1}+2u_{2}+3u_{3}+\cdots=n$
(because $\left(  u_{1},u_{2},u_{3},\ldots\right)  \in Q_{n}$). Thus, we have
defined a partition $\pi\left(  u_{1},u_{2},u_{3},\ldots\right)  $ of $n$ for
each $\left(  u_{1},u_{2},u_{3},\ldots\right)  \in Q_{n}$. In other words, we
have defined a map
\[
\pi:Q_{n}\rightarrow\left\{  \text{partitions of }n\right\}  .
\]
It remains to show that this map $\pi$ is a bijection. We define a map%
\[
\rho:\left\{  \text{partitions of }n\right\}  \rightarrow Q_{n}%
\]
that sends each partition $\lambda$ of $n$ to the sequence%
\[
\left(  \text{\# of }1\text{'s in }\lambda\text{,\ \ \# of }2\text{'s in
}\lambda\text{,\ \ \# of }3\text{'s in }\lambda\text{, \ }\ldots\right)  \in
Q_{n}%
\]
(this is indeed a sequence in $Q_{n}$, since
\[
\sum_{i=1}^{\infty}i\left(  \text{\# of }i\text{'s in }\lambda\right)
=\left(  \text{the sum of all entries of }\lambda\right)  =\left\vert
\lambda\right\vert =n
\]
because $\lambda$ is a partition of $n$). It is now easy to check that the
maps $\pi$ and $\rho$ are mutually inverse, so that $\pi$ is a bijection. The
bijection principle therefore yields $\left\vert Q_{n}\right\vert =\left(
\text{\# of partitions of }n\right)  =p\left(  n\right)  $; but this is
precisely what we wanted to show. The proof of Theorem \ref{thm.pars.main-gf}
is thus complete.
\end{proof}

Theorem \ref{thm.pars.main-gf} has a \textquotedblleft
finite\textquotedblright\ analogue (finite in the sense that the product
$\prod_{k=1}^{\infty}\dfrac{1}{1-x^{k}}$ is replaced by a finite product; the
FPSs are still infinite):

\begin{theorem}
\label{thm.pars.main-gf-parts-n}Let $m\in\mathbb{N}$. For each $n\in
\mathbb{N}$, let $p_{\operatorname*{parts}\leq m}\left(  n\right)  $ be the \#
of partitions $\lambda$ of $n$ such that all parts of $\lambda$ are $\leq m$.
Then,
\[
\sum_{n\in\mathbb{N}}p_{\operatorname*{parts}\leq m}\left(  n\right)
x^{n}=\prod_{k=1}^{m}\dfrac{1}{1-x^{k}}.
\]

\end{theorem}

\begin{proof}
[Proof of Theorem \ref{thm.pars.main-gf-parts-n}.]This proof is mostly
analogous to the above proof of Theorem \ref{thm.pars.main-gf}, and to some
extent even simpler because it uses $m$-tuples instead of infinite sequences.

We have%
\begin{align*}
&  \prod_{k=1}^{m}\underbrace{\dfrac{1}{1-x^{k}}}_{=1+x^{k}+x^{2k}%
+x^{3k}+\cdots}\\
&  =\prod_{k=1}^{m}\underbrace{\left(  1+x^{k}+x^{2k}+x^{3k}+\cdots\right)
}_{=\sum\limits_{u\in\mathbb{N}}x^{ku}}=\prod_{k=1}^{m}\ \ \sum_{u\in
\mathbb{N}}x^{ku}\\
&  =\sum_{\left(  u_{1},u_{2},\ldots,u_{m}\right)  \in\mathbb{N}^{m}}%
x^{1u_{1}}x^{2u_{2}}\cdots x^{mu_{m}}\ \ \ \ \ \ \ \ \ \ \left(
\begin{array}
[c]{c}%
\text{here, we expanded the product}\\
\text{using Proposition \ref{prop.fps.prodrule-fin-inf}}%
\end{array}
\right) \\
&  =\sum_{\left(  u_{1},u_{2},\ldots,u_{m}\right)  \in\mathbb{N}^{m}}%
x^{1u_{1}+2u_{2}+\cdots+mu_{m}}=\sum_{n\in\mathbb{N}}\left\vert Q_{n}%
\right\vert x^{n},
\end{align*}
where
\[
Q_{n}=\left\{  \left(  u_{1},u_{2},\ldots,u_{m}\right)  \in\mathbb{N}%
^{m}\text{ }\mid\ 1u_{1}+2u_{2}+\cdots+mu_{m}=n\right\}  .
\]


Thus, it will suffice to show that%
\[
\left\vert Q_{n}\right\vert =p_{\operatorname*{parts}\leq m}\left(  n\right)
\ \ \ \ \ \ \ \ \ \ \text{for each }n\in\mathbb{N}.
\]


Let us fix $n\in\mathbb{N}$. We want to construct a bijection from $Q_{n}$ to
the set $\left\{  \text{partitions }\lambda\text{ of }n\text{ such that all
parts of }\lambda\text{ are }\leq m\right\}  $.

Here is how to do this: For any $\left(  u_{1},u_{2},\ldots,u_{m}\right)  \in
Q_{n}$, define a partition%
\begin{align*}
\pi\left(  u_{1},u_{2},\ldots,u_{m}\right)  :=  &  \left(  \text{the partition
that contains each }i\text{ exactly }u_{i}\text{ times}\right) \\
=  &  \left(  \underbrace{m,m,\ldots,m}_{u_{m}\text{ times}},\ldots
,\underbrace{2,2,\ldots,2}_{u_{2}\text{ times}},\underbrace{1,1,\ldots
,1}_{u_{1}\text{ times}}\right)  .
\end{align*}
Thus, we have defined a map
\[
\pi:Q_{n}\rightarrow\left\{  \text{partitions }\lambda\text{ of }n\text{ such
that all parts of }\lambda\text{ are }\leq m\right\}  .
\]
It is easy to see that this map $\pi$ is a bijection\footnote{The argument is
analogous to the one used in the proof of Theorem \ref{thm.pars.main-gf}.}.
The bijection principle therefore yields
\[
\left\vert Q_{n}\right\vert =\left(  \text{\# of partitions }\lambda\text{ of
}n\text{ such that all parts of }\lambda\text{ are }\leq m\right)
=p_{\operatorname*{parts}\leq m}\left(  n\right)  ;
\]
but this is precisely what we wanted to show. The proof of Theorem
\ref{thm.pars.main-gf-parts-n} is thus complete.
\end{proof}

Theorem \ref{thm.pars.main-gf-parts-n} can be generalized further: Instead of
counting the partitions of $n$ whose all parts are $\leq m$, we can count the
partitions of $n$ whose parts all belong to a given set $I$ of positive
integers. This leads to the following formula (which specializes back to
Theorem \ref{thm.pars.main-gf-parts-n} when we take $I=\left\{  1,2,\ldots
,m\right\}  $):

\begin{theorem}
\label{thm.pars.main-gf-parts-I}Let $I$ be a subset of $\left\{
1,2,3,\ldots\right\}  $. For each $n\in\mathbb{N}$, let $p_{I}\left(
n\right)  $ be the \# of partitions $\lambda$ of $n$ such that all parts of
$\lambda$ belong to $I$. Then,
\[
\sum_{n\in\mathbb{N}}p_{I}\left(  n\right)  x^{n}=\prod_{k\in I}\dfrac
{1}{1-x^{k}}.
\]

\end{theorem}

\begin{proof}
[Proof of Theorem \ref{thm.pars.main-gf-parts-I} (sketched).]This is analogous
to the proof of Theorem \ref{thm.pars.main-gf-parts-n}, with some minor
changes: The $\prod_{k=1}^{m}$ sign has to be replaced by $\prod_{k\in I}$;
the $m$-tuples $\left(  u_{1},u_{2},\ldots,u_{m}\right)  \in\mathbb{N}^{m}$
must be replaced by the essentially finite families $\left(  u_{i}\right)
_{i\in I}\in\mathbb{N}^{I}$; the bijection $\pi$ has to be replaced by the new
bijection
\[
\pi:Q_{n}\rightarrow\left\{  \text{partitions }\lambda\text{ of }n\text{ such
that all parts of }\lambda\text{ are }\in I\right\}
\]
(where $Q_{n}$ is now the set of all essentially finite families $\left(
u_{i}\right)  _{i\in I}\in\mathbb{N}^{I}$) defined by%
\begin{align*}
\pi\left(  \left(  u_{i}\right)  _{i\in I}\right)  :=  &  \left(  \text{the
partition that contains each }i\in I\text{ exactly }u_{i}\text{ times}\right)
\\
=  &  \left(  \ldots,\underbrace{i_{3},i_{3},\ldots,i_{3}}_{u_{i_{3}}\text{
times}},\underbrace{i_{2},i_{2},\ldots,i_{2}}_{u_{i_{2}}\text{ times}%
},\underbrace{i_{1},i_{1},\ldots,i_{1}}_{u_{i_{1}}\text{ times}}\right)  ,
\end{align*}
where $i_{1},i_{2},i_{3},\ldots$ are the elements of $I$ listed in increasing
order. The details are left to the reader.
\end{proof}

\subsubsection{Odd parts and distinct parts}

Next, we shall state a result of Euler that we have already discovered in a
different language.

\begin{definition}
\label{def.pars.odd-dist-parts}Let $n\in\mathbb{Z}$. \medskip

\textbf{(a)} A \emph{partition of }$n$\emph{ into odd parts} means a partition
of $n$ whose all parts are odd. \medskip

\textbf{(b)} A \emph{partition of }$n$\emph{ into distinct parts} means a
partition of $n$ whose parts are distinct. \medskip

\textbf{(c)} Let%
\begin{align*}
p_{\operatorname*{odd}}\left(  n\right)   &  :=\left(  \text{\# of partitions
of }n\text{ into odd parts}\right)  \ \ \ \ \ \ \ \ \ \ \text{and}\\
p_{\operatorname*{dist}}\left(  n\right)   &  :=\left(  \text{\# of partitions
of }n\text{ into distinct parts}\right)  .
\end{align*}

\end{definition}

\begin{example}
We have%
\begin{align*}
p_{\operatorname*{odd}}\left(  7\right)   &  =\left\vert \left\{  \left(
7\right)  ,\ \ \left(  5,1,1\right)  ,\ \ \left(  3,3,1\right)  ,\ \ \left(
3,1,1,1,1\right)  ,\ \ \left(  1,1,1,1,1,1,1\right)  \right\}  \right\vert
=5;\\
p_{\operatorname*{dist}}\left(  7\right)   &  =\left\vert \left\{  \left(
7\right)  ,\ \ \left(  6,1\right)  ,\ \ \left(  5,2\right)  ,\ \ \left(
4,3\right)  ,\ \ \left(  4,2,1\right)  \right\}  \right\vert =5.
\end{align*}

\end{example}

\begin{theorem}
[Euler's odd-distinct identity]\label{thm.pars.odd-dist-equal}We have
$p_{\operatorname*{odd}}\left(  n\right)  =p_{\operatorname*{dist}}\left(
n\right)  $ for each $n\in\mathbb{N}$.
\end{theorem}

We have already encountered this theorem before (as Theorem
\ref{thm.gf.prod.euler-comb}, albeit in less precise language), and we have
proved it using the generating function identity%
\[
\prod_{i>0}\left(  1-x^{2i-1}\right)  ^{-1}=\prod_{k>0}\left(  1+x^{k}\right)
.
\]


Let me outline a different, bijective proof.

\begin{proof}
[Second proof of Theorem \ref{thm.pars.odd-dist-equal} (sketched).]Let
$n\in\mathbb{N}$. We want to construct a bijection%
\[
A:\left\{  \text{partitions of }n\text{ into odd parts}\right\}
\rightarrow\left\{  \text{partitions of }n\text{ into distinct parts}\right\}
.
\]
We shall do this as follows: Given a partition $\lambda$ of $n$ into odd
parts, we repeatedly merge pairs of equal parts in $\lambda$ until no more
equal parts appear. The final result will be $A\left(  \lambda\right)  $. Here
are two examples:

\begin{itemize}
\item To compute $A\left(  5,5,3,1,1,1\right)  $, we compute\footnote{The two
entries underlined are the two equal entries that are going to get merged in
the next step. Note that there are usually several candidates, and we just
pick one pair at will.}%
\[
\left(  \underline{5,5},3,1,1,1\right)  \rightarrow\left(
10,3,1,\underline{1,1}\right)  \rightarrow\left(  10,3,2,1\right)  .
\]
Thus, $A\left(  5,5,3,1,1,1\right)  =\left(  10,3,2,1\right)  $.

\item To compute $A\left(  5,3,1,1,1,1\right)  $, we compute%
\[
\left(  5,3,1,1,\underline{1,1}\right)  \rightarrow\left(
5,3,2,\underline{1,1}\right)  \rightarrow\left(  5,3,\underline{2,2}\right)
\rightarrow\left(  5,4,3\right)  .
\]
Thus, $A\left(  5,3,1,1,1,1\right)  =\left(  5,4,3\right)  $.
\end{itemize}

Why is this map $A$ well-defined? We only specified the sort of steps we are
allowed to take when computing $A\left(  \lambda\right)  $; however, there is
often a choice involved in taking these steps (since there are often several
pairs of equal parts).\footnote{For example, we could have also computed
$A\left(  5,5,3,1,1,1\right)  $ as follows:%
\[
\left(  5,5,3,\underline{1},1,\underline{1}\right)  \rightarrow\left(
\underline{5,5},3,2,1\right)  \rightarrow\left(  10,3,2,1\right)  .
\]
} So we have specified a non-deterministic algorithm. Why is the resulting
partition independent of the choices we make?

One way to prove this is using the \emph{diamond lemma}, which is a general
tool for proving that certain non-deterministic algorithms have unique final
outcomes (independent of the choices taken). See, e.g.,
\url{https://mathoverflow.net/questions/289300/} for a list of references on
this lemma.

For the map $A$, we can also proceed differently, by analyzing the algorithm
that we used to define $A$. Namely, we observe what is really going on when we
are merging equal parts. Let us say our original partition $\lambda$ has $p$
many $1$s. Let us first merge them in pairs, so that we get $\left\lfloor
p/2\right\rfloor $ many $2$s and maybe one single $1$. Then, let us merge the
$2$s in pairs, so that we get $\left\lfloor \left\lfloor p/2\right\rfloor
/2\right\rfloor $ many $4$s, maybe a single $2$, and maybe a single $1$.
Proceed until no more than one $1$, no more than one $2$, no more than one
$4$, no more than one $8$, and so on remain. This clears out any duplicate
parts of the form $2^{k}$. Next do the same with parts of the form
$3\cdot2^{k}$ (that is, with parts equal to $3$, $6$, $12$, $24$, and so on),
then with parts of the form $5\cdot2^{k}$, and so on.

The nice thing about this way of proceeding is that we can explicitly describe
the final outcome. Indeed, if the original partition $\lambda$ (a partition of
$n$ into odd parts) contains an odd part $k$ precisely $m$ many times, and if
the binary representation of $m$ is $m=\left(  m_{i}m_{i-1}\cdots m_{1}%
m_{0}\right)  _{2}$ (that is, if $m_{0},m_{1},\ldots,m_{i}\in\left\{
0,1\right\}  $ satisfy $m=\sum_{j=0}^{i}m_{j}2^{j}$), then the partition
$A\left(  \lambda\right)  $ will contain the number $2^{0}k$ exactly $m_{0}$
times, the number $2^{1}k$ exactly $m_{1}$ times, the number $2^{2}k$ exactly
$m_{2}$ times, and so on. Since the binary digits $m_{0},m_{1},\ldots,m_{i}$
are all $\leq1$, this partition $A\left(  \lambda\right)  $ will therefore not
contain any number more than once, i.e., it will be a partition into distinct parts.

It is not hard to check that this map $A$ is indeed a bijection. Indeed, in
order to see this, we construct a map $B$ that will turn out to be its
inverse. Here, we start with a partition $\lambda$ of $n$ into distinct parts.
Let us represent each part of this partition in the form $k\cdot2^{i}$ for
some odd $k\geq1$ and some integer $i\geq0$. (Recall that any positive integer
can be represented uniquely in this form.) Now, replace this part $k\cdot
2^{i}$ by $2^{i}$ many $k$'s. The resulting partition (once all parts have
been replaced) will usually have many equal parts, but all its parts are odd.
We define $B\left(  \lambda\right)  $ to be this resulting partition.
Alternatively, $B\left(  \lambda\right)  $ can also be constructed
step-by-step by a non-deterministic algorithm: Starting with $\lambda$, keep
\textquotedblleft breaking even parts into halves\textquotedblright\ (i.e.,
whenever you see an even part $m$, replace it by two parts $\dfrac{m}{2}$ and
$\dfrac{m}{2}$), until no even parts remain any more. The result is $B\left(
\lambda\right)  $. It is not hard to see that both descriptions of $B\left(
\lambda\right)  $ describe the same partition. It is furthermore easy to see
that this map $B$ is indeed an inverse of $A$, so that $A$ is indeed a
bijection. Thus, the bijection principle yields%
\[
\left\vert \left\{  \text{partitions of }n\text{ into odd parts}\right\}
\right\vert =\left\vert \left\{  \text{partitions of }n\text{ into distinct
parts}\right\}  \right\vert .
\]
In other words, $p_{\operatorname*{odd}}\left(  n\right)
=p_{\operatorname*{dist}}\left(  n\right)  $. This proves Theorem
\ref{thm.pars.odd-dist-equal}.
\end{proof}

Yet another proof of Theorem \ref{thm.pars.odd-dist-equal} will be given in
Subsection \ref{subsec.sign.pie.exas}.

\subsubsection{Partitions with a given largest part}

Here is another situation in which two kinds of partitions are equinumerous:

\begin{proposition}
\label{prop.pars.pkn=dual}Let $n\in\mathbb{N}$ and $k\in\mathbb{N}$. Then,%
\[
p_{k}\left(  n\right)  =\left(  \text{\# of partitions of }n\text{ whose
largest part is }k\right)  .
\]

\end{proposition}

Here and in the following, we use the following convention:

\begin{convention}
\label{conv.pars.largest-part-0}We agree to say that the largest part of the
empty partition $\left(  {}\right)  $ is $0$ (even though this partition has
no parts).
\end{convention}

\begin{example}
For $n=4$ and $k=3$, we have%
\[
p_{k}\left(  n\right)  =p_{3}\left(  4\right)  =1\ \ \ \ \ \ \ \ \ \ \left(
\text{due to the partition }\left(  2,1,1\right)  \right)
\]
and%
\begin{align*}
&  \left(  \text{\# of partitions of }n\text{ whose largest part is }k\right)
\\
&  =\left(  \text{\# of partitions of }4\text{ whose largest part is }3\right)
\\
&  =1\ \ \ \ \ \ \ \ \ \ \left(  \text{due to the partition }\left(
3,1\right)  \right)  .
\end{align*}
Thus, Proposition \ref{prop.pars.pkn=dual} holds for $n=4$ and $k=3$.
\end{example}

\begin{proof}
[Proof of Proposition \ref{prop.pars.pkn=dual} (sketched).]We do a
\textquotedblleft proof by picture\textquotedblright\ (it can be made rigorous
-- see Exercise \ref{exe.pars.transpose} for this). We pick $n=14$ and $k=4$
for example, and we start with the partition $\lambda=\left(  5,4,4,1\right)
$ of $n$ into $k$ parts.

We draw a table of $k$ left-aligned rows, where the length of each row equals
the corresponding part of $\lambda$ (that is, the $i$-th row from the top has
$\lambda_{i}$ boxes, where $\lambda=\left(  \lambda_{1},\lambda_{2}%
,\ldots,\lambda_{k}\right)  $):%
\[%
% [tableau figure removed]
%BeginExpansion
\begin{ytableau}
\none[5 \to] & \none& & & & & \\
\none[4 \to] & \none& & & & \\
\none[4 \to] & \none& & & & \\
\none[1 \to] & \none&
\end{ytableau}%
%EndExpansion
\]


Now, let us flip this table across the \textquotedblleft main
diagonal\textquotedblright\ (i.e., the diagonal that goes from the top-left
corner to the bottom-right corner)\footnote{This kind of flip is precisely how
you would transpose a matrix.}:%
\[%
% [tableau figure removed]
%BeginExpansion
\begin{ytableau}
\none[5] & \none[4] & \none[4] & \none[1] \\
\none[\downarrow] & \none[\downarrow] & \none[\downarrow] & \none
[\downarrow] \\
& & & \\
& & \\
& & \\
& & \\
\;
\end{ytableau}%
%EndExpansion
\]


The lengths of the rows of the resulting table again form a partition of $n$.
(In our case, this new partition is $\left(  4,3,3,3,1\right)  $.) Moreover,
the largest part of this new partition is $k$ (because the original table had
$k$ rows, so the flipped table has $k$ columns, and this means that its top
row has $k$ boxes). This procedure (i.e., turning a partition into a table,
then flipping the table across the \textquotedblleft main
diagonal\textquotedblright, and then reading the lengths of the rows of the
resulting table again as a partition) therefore gives a map from%
\[
\left\{  \text{partitions of }n\text{ into }k\text{ parts}\right\}
\]
to%
\[
\left\{  \text{partitions of }n\text{ whose largest part is }k\right\}  .
\]
Moreover, this map is a bijection (indeed, its inverse can be effected in the
exact same way, by flipping the table). This bijection is called
\emph{conjugation} of partitions, and will be studied in more detail later.

Here are some pointers to how this proof can be formalized (see Exercise
\ref{exe.pars.transpose} for much more): For any partition $\lambda=\left(
\lambda_{1},\lambda_{2},\ldots,\lambda_{k}\right)  $, we define the
\emph{Young diagram} of $\lambda$ to be the set%
\[
Y\left(  \lambda\right)  :=\left\{  \left(  i,j\right)  \in\mathbb{Z}%
^{2}\ \mid\ 1\leq i\leq k\text{ and }1\leq j\leq\lambda_{i}\right\}  .
\]
This Young diagram is precisely the table that we drew above, as long as we
agree to identify each pair $\left(  i,j\right)  \in Y\left(  \lambda\right)
$ with the box in row $i$ and column $j$. Now, the \emph{conjugate} of the
partition $\lambda$ is the partition $\lambda^{t}$ uniquely determined by%
\[
Y\left(  \lambda^{t}\right)  =\operatorname*{flip}\left(  Y\left(
\lambda\right)  \right)  =\left\{  \left(  j,i\right)  \ \mid\ \left(
i,j\right)  \in Y\left(  \lambda\right)  \right\}  .
\]
Explicitly, $\lambda^{t}$ can be defined by $\lambda^{t}=\left(  \mu_{1}%
,\mu_{2},\ldots,\mu_{p}\right)  $, where $p$ is the largest part of $\lambda$
and where%
\[
\mu_{i}=\left(  \text{\# of parts of }\lambda\text{ that are }\geq i\right)
\ \ \ \ \ \ \ \ \ \ \text{for each }i\in\left\{  1,2,\ldots,p\right\}  .
\]
(This conjugate $\lambda^{t}$ is also often called $\lambda^{\prime}$, and is
also known as the \emph{transpose} of $\lambda$.) Now, it is not hard to show
that $\left\vert \lambda^{t}\right\vert =\left\vert \lambda\right\vert $ and
$\left(  \lambda^{t}\right)  ^{t}=\lambda$ for each partition $\lambda$, and
that the largest part of $\lambda^{t}$ equals the length of $\lambda$. Using
these observations (which are proved in Exercise \ref{exe.pars.transpose}), we
see that the map%
\begin{align*}
\left\{  \text{partitions of }n\text{ into }k\text{ parts}\right\}   &
\rightarrow\left\{  \text{partitions of }n\text{ whose largest part is
}k\right\}  ,\\
\lambda &  \mapsto\lambda^{t}%
\end{align*}
is well-defined and is a bijection; thus, the above proof of Proposition
\ref{prop.pars.pkn=dual} becomes fully rigorous.
\end{proof}

The word \textquotedblleft Young\textquotedblright\ in \textquotedblleft Young
diagram\textquotedblright\ (and, later, \textquotedblleft Young
tableau\textquotedblright) does not imply any novelty (Young diagrams have
been around in some form or another since the 19th century -- if often in the
superficially different guise of \textquotedblleft Ferrers
diagrams\textquotedblright), but rather honors
\href{https://en.wikipedia.org/wiki/Alfred_Young}{Alfred Young}, who built up
the representation theory of symmetric groups (and significantly forwarded
invariant theory) using these objects.

\begin{corollary}
\label{cor.pars.p0kn=dual}Let $n\in\mathbb{N}$ and $k\in\mathbb{N}$. Then,%
\begin{align*}
&  p_{0}\left(  n\right)  +p_{1}\left(  n\right)  +\cdots+p_{k}\left(
n\right) \\
&  =\left(  \text{\# of partitions of }n\text{ whose largest part is }\leq
k\right)  .
\end{align*}

\end{corollary}

\begin{proof}
[Proof of Corollary \ref{cor.pars.p0kn=dual}.]We have%
\begin{align*}
&  p_{0}\left(  n\right)  +p_{1}\left(  n\right)  +\cdots+p_{k}\left(
n\right) \\
&  =\sum_{i=0}^{k}\underbrace{p_{i}\left(  n\right)  }_{\substack{=\left(
\text{\# of partitions of }n\text{ whose largest part is }i\right)
\\\text{(by Proposition \ref{prop.pars.pkn=dual},}\\\text{applied to }i\text{
instead of }k\text{)}}}\\
&  =\sum_{i=0}^{k}\left(  \text{\# of partitions of }n\text{ whose largest
part is }i\right) \\
&  =\left(  \text{\# of partitions of }n\text{ whose largest part is }\leq
k\right)  .
\end{align*}
This proves Corollary \ref{cor.pars.p0kn=dual}.
\end{proof}

Corollary \ref{cor.pars.p0kn=dual} leads to yet another FPS identity:

\begin{theorem}
\label{thm.pars.main-gf-0n}Let $m\in\mathbb{N}$. Then,%
\[
\sum_{n\in\mathbb{N}}\left(  p_{0}\left(  n\right)  +p_{1}\left(  n\right)
+\cdots+p_{m}\left(  n\right)  \right)  x^{n}=\prod_{k=1}^{m}\dfrac{1}%
{1-x^{k}}.
\]

\end{theorem}

\begin{proof}
[Proof of Theorem \ref{thm.pars.main-gf-0n}.]For each $n\in\mathbb{N}$, we
have%
\begin{align*}
&  p_{0}\left(  n\right)  +p_{1}\left(  n\right)  +\cdots+p_{m}\left(
n\right) \\
&  =\left(  \text{\# of partitions of }n\text{ whose largest part is }\leq
m\right)  \ \ \ \ \ \ \ \ \ \ \left(  \text{by Corollary
\ref{cor.pars.p0kn=dual}}\right) \\
&  =\left(  \text{\# of partitions of }n\text{ whose all parts are }\leq
m\right) \\
&  \ \ \ \ \ \ \ \ \ \ \ \ \ \ \ \ \ \ \ \ \left(
\begin{array}
[c]{c}%
\text{because the condition \textquotedblleft the largest part is }\leq
m\text{\textquotedblright\ for a}\\
\text{partition is clearly equivalent to \textquotedblleft all parts are }\leq
m\text{\textquotedblright}%
\end{array}
\right) \\
&  =p_{\operatorname*{parts}\leq m}\left(  n\right)  ,
\end{align*}
where $p_{\operatorname*{parts}\leq m}\left(  n\right)  $ is defined as in
Theorem \ref{thm.pars.main-gf-parts-n}. Hence,%
\[
\sum_{n\in\mathbb{N}}\underbrace{\left(  p_{0}\left(  n\right)  +p_{1}\left(
n\right)  +\cdots+p_{m}\left(  n\right)  \right)  }_{=p_{\operatorname*{parts}%
\leq m}\left(  n\right)  }x^{n}=\sum_{n\in\mathbb{N}}p_{\operatorname*{parts}%
\leq m}\left(  n\right)  x^{n}=\prod_{k=1}^{m}\dfrac{1}{1-x^{k}}%
\]
(by Theorem \ref{thm.pars.main-gf-parts-n}). This proves Theorem
\ref{thm.pars.main-gf-0n}.
\end{proof}

\subsubsection{Partition number vs. sums of divisors}

We shall now prove a curious combinatorial identity (due, I think, to Euler),
illustrating the usefulness of generating functions and logarithmic
derivatives. It connects the partition numbers $p\left(  n\right)  $ with a
further sequence, which counts the positive divisors of a given positive integer:

\begin{theorem}
\label{thm.pars.sigma1}For any positive integer $n$, let $\sigma\left(
n\right)  $ denote the sum of all positive divisors of $n$. (For example,
$\sigma\left(  6\right)  =1+2+3+6=12$ and $\sigma\left(  7\right)  =1+7=8$.)

For any $n\in\mathbb{N}$, we have%
\[
np\left(  n\right)  =\sum_{k=1}^{n}\sigma\left(  k\right)  p\left(
n-k\right)  .
\]

\end{theorem}

For example, for $n=3$, this says that%
\[
3\underbrace{p\left(  3\right)  }_{=3}=\underbrace{\sigma\left(  1\right)
}_{=1}\underbrace{p\left(  2\right)  }_{=2}+\underbrace{\sigma\left(
2\right)  }_{=3}\underbrace{p\left(  1\right)  }_{=1}+\underbrace{\sigma
\left(  3\right)  }_{=4}\underbrace{p\left(  0\right)  }_{=1}.
\]


For reference, let us give a table of the first $15$ sums $\sigma\left(
n\right)  $ defined in Theorem \ref{thm.pars.sigma1} (see the
\href{https://oeis.org/A000203}{Sequence A000203 in the OEIS} for more
values):%
\[%
\begin{tabular}
[c]{|c||c|c|c|c|c|c|c|c|c|c|c|c|c|c|c|}\hline
$n$ & $1$ & $2$ & $3$ & $4$ & $5$ & $6$ & $7$ & $8$ & $9$ & $10$ & $11$ & $12$
& $13$ & \multicolumn{1}{|c|}{$14$} & $15$\\\hline
$\sigma\left(  n\right)  $ & $1$ & $3$ & $4$ & $7$ & $6$ & $12$ & $8$ & $15$ &
$13$ & $18$ & $12$ & $28$ & $14$ & $24$ & $24$\\\hline
\end{tabular}
\ \ \ \ \ .
\]


\begin{proof}
[Proof of Theorem \ref{thm.pars.sigma1} (sketched).]Define the two FPSs%
\[
P:=\sum_{n\in\mathbb{N}}p\left(  n\right)  x^{n}\in\mathbb{Z}\left[  \left[
x\right]  \right]  \ \ \ \ \ \ \ \ \ \ \text{and}\ \ \ \ \ \ \ \ \ \ S:=\sum
_{k>0}\sigma\left(  k\right)  x^{k}\in\mathbb{Z}\left[  \left[  x\right]
\right]  .
\]
Hence,%
\begin{align}
SP  &  =\left(  \sum_{k>0}\sigma\left(  k\right)  x^{k}\right)  \left(
\sum_{n\in\mathbb{N}}p\left(  n\right)  x^{n}\right)  =\sum_{k>0}%
\ \ \sum_{n\in\mathbb{N}}\sigma\left(  k\right)  \underbrace{x^{k}p\left(
n\right)  x^{n}}_{=p\left(  n\right)  x^{k+n}}\nonumber\\
&  =\sum_{k>0}\ \ \sum_{n\in\mathbb{N}}\sigma\left(  k\right)  p\left(
n\right)  x^{k+n}=\underbrace{\sum_{k>0}\ \ \sum_{\substack{m\in
\mathbb{N};\\m\geq k}}}_{\substack{=\sum_{m\in\mathbb{N}}\ \ \sum
_{\substack{k>0;\\m\geq k}}\\=\sum_{m\in\mathbb{N}}\ \ \sum_{k=1}^{m}}%
}\sigma\left(  k\right)  p\left(  m-k\right)  \underbrace{x^{k+\left(
m-k\right)  }}_{=x^{m}}\nonumber\\
&  \ \ \ \ \ \ \ \ \ \ \ \ \ \ \ \ \ \ \ \ \left(  \text{here, we have
substituted }m-k\text{ for }n\text{ in the inner sum}\right) \nonumber\\
&  =\sum_{m\in\mathbb{N}}\ \ \sum_{k=1}^{m}\sigma\left(  k\right)  p\left(
m-k\right)  x^{m}. \label{pf.thm.pars.sigma1.SP=}%
\end{align}
From $P=\sum_{n\in\mathbb{N}}p\left(  n\right)  x^{n}$, we obtain%
\begin{equation}
P^{\prime}=\sum_{n>0}np\left(  n\right)  x^{n-1}.
\label{pf.thm.pars.sigma1.P'=}%
\end{equation}
Multiplying this equality by $x$, we obtain the somewhat nicer%
\begin{align}
xP^{\prime}  &  =x\sum_{n>0}np\left(  n\right)  x^{n-1}=\sum_{n>0}np\left(
n\right)  \underbrace{xx^{n-1}}_{=x^{n}}=\sum_{n>0}np\left(  n\right)
x^{n}\nonumber\\
&  =\sum_{n\in\mathbb{N}}np\left(  n\right)  x^{n}
\label{pf.thm.pars.sigma1.xP'=}%
\end{align}
(here, we have extended the range of the sum to include $n=0$, which caused no
change to the value of the sum, since the newly added $n=0$ addend is
$0p\left(  0\right)  x^{0}=0$).

Let us now recall the notion of a logarithmic derivative (as defined in
Definition \ref{def.fps.loder.1}). The FPS $P$ has constant coefficient
$\left[  x^{0}\right]  P=p\left(  0\right)  =1$, thus belongs to
$\mathbb{Z}\left[  \left[  x\right]  \right]  _{1}$. Hence, its logarithmic
derivative $\operatorname*{loder}P$ is well-defined.

We shall now use $\operatorname*{loder}P$ to compute $P^{\prime}$ in a
roundabout way. Namely, we have%
\[
P=\sum_{n\in\mathbb{N}}p\left(  n\right)  x^{n}=\prod_{k=1}^{\infty}\dfrac
{1}{1-x^{k}}%
\]
(by Theorem \ref{thm.pars.main-gf}). In other words,%
\begin{equation}
P=\prod_{k>0}\dfrac{1}{1-x^{k}} \label{pf.thm.pars.sigma1.P=prod}%
\end{equation}
(since the product sign $\prod\limits_{k=1}^{\infty}$ is synonymous with
$\prod\limits_{k>0}$).

But Corollary \ref{cor.fps.loder.prodk} says that any $k$ FPSs $f_{1}%
,f_{2},\ldots,f_{k}\in K\left[  \left[  x\right]  \right]  _{1}$ (for any
commutative ring $K$) satisfy%
\[
\operatorname*{loder}\left(  f_{1}f_{2}\cdots f_{k}\right)
=\operatorname*{loder}\left(  f_{1}\right)  +\operatorname*{loder}\left(
f_{2}\right)  +\cdots+\operatorname*{loder}\left(  f_{k}\right)  .
\]
In other words, the logarithmic derivative of a product of $k$ FPSs (which
belong to $K\left[  \left[  x\right]  \right]  _{1}$) equals the sum of the
logarithmic derivatives of these $k$ FPSs. By a simple limit argument (using
Theorem \ref{thm.fps.lim.prod-lim}, Theorem \ref{thm.fps.lim.sum-lim},
Proposition \ref{prop.fps.lim.sum-quot} and Proposition
\ref{prop.fps.lim.deriv-lim}), we can extend this fact to an infinite product
(as long as it is multipliable). Thus, if $\left(  f_{1},f_{2},f_{3}%
,\ldots\right)  $ is a multipliable sequence of FPSs in $K\left[  \left[
x\right]  \right]  _{1}$ (for any commutative ring $K$), then%
\[
\operatorname*{loder}\left(  f_{1}f_{2}f_{3}\cdots\right)
=\operatorname*{loder}\left(  f_{1}\right)  +\operatorname*{loder}\left(
f_{2}\right)  +\operatorname*{loder}\left(  f_{3}\right)  +\cdots;
\]
equivalently,%
\[
\operatorname*{loder}\left(  \prod_{k>0}f_{k}\right)  =\sum_{k>0}%
\operatorname*{loder}\left(  f_{k}\right)  .
\]
Applying this to $K=\mathbb{Z}$ and $f_{k}=\dfrac{1}{1-x^{k}}$, we obtain%
\[
\operatorname*{loder}\left(  \prod_{k>0}\dfrac{1}{1-x^{k}}\right)  =\sum
_{k>0}\operatorname*{loder}\dfrac{1}{1-x^{k}}.
\]
In view of (\ref{pf.thm.pars.sigma1.P=prod}), we can rewrite this as%
\begin{equation}
\operatorname*{loder}P=\sum_{k>0}\operatorname*{loder}\dfrac{1}{1-x^{k}}.
\label{pf.thm.pars.sigma1.loderP=sumk}%
\end{equation}
However, every positive integer $k$ satisfies%
\begin{align*}
\operatorname*{loder}\underbrace{\dfrac{1}{1-x^{k}}}_{=\left(  1-x^{k}\right)
^{-1}}  &  =\operatorname*{loder}\left(  \left(  1-x^{k}\right)  ^{-1}\right)
\\
&  =-\underbrace{\operatorname*{loder}\left(  1-x^{k}\right)  }%
_{\substack{=\dfrac{\left(  1-x^{k}\right)  ^{\prime}}{1-x^{k}}\\\text{(by the
definition of a}\\\text{logarithmic derivative)}}}\ \ \ \ \ \ \ \ \ \ \left(
\begin{array}
[c]{c}%
\text{by Corollary \ref{cor.fps.loder.inv},}\\
\text{applied to }f=1-x^{k}%
\end{array}
\right) \\
&  =-\dfrac{\left(  1-x^{k}\right)  ^{\prime}}{1-x^{k}}=-\dfrac{-kx^{k-1}%
}{1-x^{k}}\ \ \ \ \ \ \ \ \ \ \left(  \text{since }\left(  1-x^{k}\right)
^{\prime}=-kx^{k-1}\right) \\
&  =\dfrac{kx^{k-1}}{1-x^{k}}=kx^{k-1}\underbrace{\left(  1-x^{k}\right)
^{-1}}_{\substack{=\sum_{m\in\mathbb{N}}\left(  x^{k}\right)  ^{m}\\\text{(by
the geometric series formula)}}}\\
&  =kx^{k-1}\sum_{m\in\mathbb{N}}\left(  x^{k}\right)  ^{m}=\sum
_{m\in\mathbb{N}}k\underbrace{x^{k-1}\left(  x^{k}\right)  ^{m}}%
_{=x^{k-1+km}=x^{k\left(  m+1\right)  -1}}\\
&  =\sum_{m\in\mathbb{N}}kx^{k\left(  m+1\right)  -1}\\
&  =\sum_{i>0}kx^{ki-1}\ \ \ \ \ \ \ \ \ \ \left(
\begin{array}
[c]{c}%
\text{here, we have substituted }i\\
\text{for }m+1\text{ in the sum}%
\end{array}
\right) \\
&  =\sum_{\substack{n>0;\\k\mid n}}kx^{n-1}\ \ \ \ \ \ \ \ \ \ \left(
\begin{array}
[c]{c}%
\text{here, we have substituted }n\\
\text{for }ki\text{ in the sum}%
\end{array}
\right)  .
\end{align*}
Hence, we can rewrite (\ref{pf.thm.pars.sigma1.loderP=sumk}) as%
\begin{align*}
\operatorname*{loder}P  &  =\underbrace{\sum_{k>0}\ \ \sum
_{\substack{n>0;\\k\mid n}}}_{=\sum_{n>0}\ \ \sum_{\substack{k>0;\\k\mid n}%
}}kx^{n-1}=\sum_{n>0}\ \ \sum_{\substack{k>0;\\k\mid n}}kx^{n-1}\\
&  =\sum_{n>0}\underbrace{\left(  \sum_{\substack{k>0;\\k\mid n}}k\right)
}_{\substack{=\left(  \text{sum of all positive divisors of }n\right)
\\=\sigma\left(  n\right)  \\\text{(by the definition of }\sigma\left(
n\right)  \text{)}}}x^{n-1}=\sum_{n>0}\sigma\left(  n\right)  x^{n-1}.
\end{align*}
Comparing this with%
\[
\operatorname*{loder}P=\dfrac{P^{\prime}}{P}\ \ \ \ \ \ \ \ \ \ \left(
\text{by Definition \ref{def.fps.loder.1}}\right)  ,
\]
we obtain
\[
\dfrac{P^{\prime}}{P}=\sum_{n>0}\sigma\left(  n\right)  x^{n-1}.
\]
Multiplying both sides of this equality by $xP$, we find%
\begin{align*}
xP^{\prime}  &  =xP\cdot\sum_{n>0}\sigma\left(  n\right)  x^{n-1}=P\cdot
\sum_{n>0}\sigma\left(  n\right)  \underbrace{xx^{n-1}}_{=x^{n}}%
=P\cdot\underbrace{\sum_{n>0}\sigma\left(  n\right)  x^{n}}_{=\sum_{k>0}%
\sigma\left(  k\right)  x^{k}=S}=P\cdot S=SP\\
&  =\sum_{m\in\mathbb{N}}\ \ \sum_{k=1}^{m}\sigma\left(  k\right)  p\left(
m-k\right)  x^{m}\ \ \ \ \ \ \ \ \ \ \left(  \text{by
(\ref{pf.thm.pars.sigma1.SP=})}\right) \\
&  =\sum_{n\in\mathbb{N}}\ \ \sum_{k=1}^{n}\sigma\left(  k\right)  p\left(
n-k\right)  x^{n}\ \ \ \ \ \ \ \ \ \ \left(
\begin{array}
[c]{c}%
\text{here, we have renamed the}\\
\text{summation index }m\text{ as }n
\end{array}
\right)  .
\end{align*}
Comparing this with (\ref{pf.thm.pars.sigma1.xP'=}), we find%
\[
\sum_{n\in\mathbb{N}}np\left(  n\right)  x^{n}=\sum_{n\in\mathbb{N}}%
\ \ \sum_{k=1}^{n}\sigma\left(  k\right)  p\left(  n-k\right)  x^{n}.
\]
Comparing coefficients in front of $x^{n}$ on both sides of this equality, we
find that%
\[
np\left(  n\right)  =\sum_{k=1}^{n}\sigma\left(  k\right)  p\left(
n-k\right)  \ \ \ \ \ \ \ \ \ \ \text{for each }n\in\mathbb{N}\text{.}%
\]
This proves Theorem \ref{thm.pars.sigma1}.
\end{proof}

More generally, the following holds:

\begin{theorem}
\label{thm.pars.sigma1-I}Let $I$ be a subset of $\left\{  1,2,3,\ldots
\right\}  $. For each $n\in\mathbb{N}$, let $p_{I}\left(  n\right)  $ be the
\# of partitions $\lambda$ of $n$ such that all parts of $\lambda$ belong to
$I$.

For any positive integer $n$, let $\sigma_{I}\left(  n\right)  $ denote the
sum of all positive divisors of $n$ that belong to $I$. (For example, if
$O=\left\{  \text{all odd positive integers}\right\}  $ and $E=\left\{
\text{all even positive integers}\right\}  $, then $\sigma_{O}\left(
6\right)  =1+3=4$ and $\sigma_{E}\left(  6\right)  =2+6=8$.)

For any $n\in\mathbb{N}$, we have%
\[
np_{I}\left(  n\right)  =\sum_{k=1}^{n}\sigma_{I}\left(  k\right)
p_{I}\left(  n-k\right)  .
\]

\end{theorem}

Of course, Theorem \ref{thm.pars.sigma1} is the particular case of Theorem
\ref{thm.pars.sigma1-I} for $I=\left\{  1,2,3,\ldots\right\}  $.

\begin{proof}
[Proof of Theorem \ref{thm.pars.sigma1-I} (sketched).]Argue as in our above
proof of Theorem \ref{thm.pars.sigma1}, making some replacements: In
particular, symbols like $\sum\limits_{k>0}$ and $\prod\limits_{k>0}$ must be
replaced by $\sum\limits_{k\in I}$ and $\prod\limits_{k\in I}$ (respectively),
and Theorem \ref{thm.pars.main-gf-parts-I} must be used instead of Theorem
\ref{thm.pars.main-gf}. We leave the details to the reader.
\end{proof}

We note that Theorem \ref{thm.pars.sigma1} and even the more general Theorem
\ref{thm.pars.sigma1-I} are not hard to prove by completely elementary ways
(see, e.g., Exercise \ref{exe.pars.sigma-I-elem}). Nevertheless, our above
proof using generating functions has some interesting consequences, one of
which we might soon explore.
